\par\medskip\noindent\textbf{Cómo leer la salida.}
{\small\setlength{\tabcolsep}{4pt}\renewcommand{\arraystretch}{1.15}
\begin{xltabular}{\linewidth}{@{}>{\raggedright\arraybackslash}p{0.24\linewidth}>{\raggedleft\arraybackslash}p{0.15\linewidth}>{\raggedright\arraybackslash}X@{}}
\toprule Indicador & Valor & Interpretación \\ \midrule \endfirsthead
\toprule Indicador & Valor & Interpretación \\ \midrule \endhead
\bottomrule \endlastfoot
$T$ & 132 meses & Rendimientos simples mensuales de 2015 a 2025. \\
\texttt{media} & 0,0113 / 0,0101 & Rendimiento medio mensual: 1,13\,\% (\texttt{EPU}) y 1,01\,\% (S\&P 500). \\
\texttt{desv} & 0,0712 / 0,0430 & Volatilidad mensual; \texttt{EPU} es mucho más riesgoso. \\
\texttt{asimetría}, \texttt{curtosis} & $-0{,}20$ y 3,92; $-0{,}40$ y 3,52 & Desviaciones moderadas de la normalidad (asimetría 0 y curtosis 3). \\
\texttt{theta mensual} & 0,159 / 0,234 & Razón media/desviación mensual estimada, $\hat\theta=\hat\mu/\hat\sigma$. \\
Anualizada $\times\sqrt{12}$ & 0,550 / 0,811 & Escala anual habitual de la razón de Sharpe (con tasa libre de riesgo nula). \\
EE (normal) & 0,0876 / 0,0882 & Error estándar con $v=1+\theta^2/2$: muy cercano a $1/\sqrt T$. \\
IC 95\,\% (normal) & $[-0{,}013;\ 0{,}330]$ / $[0{,}061;\ 0{,}407]$ & Para \texttt{EPU} incluye el cero: no se puede afirmar que su media fue positiva. Para el S\&P 500, sí. \\
EE (general) & 0,0892 / 0,0924 & Error con asimetría y curtosis: algo mayor porque la asimetría es negativa y la curtosis supera 3. \\
IC 95\,\% (general) & $[-0{,}016;\ 0{,}334]$ / $[0{,}053;\ 0{,}415]$ & Las conclusiones no cambian. \\
\end{xltabular}}
